Program 09:
Develop a LaTeX script to create a document that consists of two paragraphs
with a minimum of 10 citations in it and display the reference in the section
\documentclass{article}
\usepackage{fancyhdr}
\pagestyle{fancy}
\fancyhead[l]{Latex,program 9}
\fancyhead[r]{by Sagar}
\fancyfoot[l]{Dept of CSE,GEC,Karwar.}
\fancyfoot[r]{2025-26}
\renewcommand{\headrulewidth}{2pt}
\renewcommand{\footrulewidth}{2pt}
\begin{document}
\section{Introduction}
Lorem ipsum dolor sit amet, consectetur adipiscing elit. Fusce condimentum nulla a nisl
consectetur, in vehicula purus fermentum. Pellentesque habitant morbi tristique senectus et
netus et malesuada fames ac turpis egestas. Sed efficitur placerat sapien ut placerat. Integer
fringilla justo ac massa ultrices, sed laoreet est fermentum.Pellentesque habitant morbi
tristique senectus et netus et malesuada fames ac turpis egestas. Donec fermentum vehicula
felis, in congue odio volutpat vitae \cite{citation1,
citation2, citation3, citation4, citation5, citation6, citation7, citation8,
citation9,citation10}.Morbi gravida mi sit amet justo aliquet, eu mollis velit convallis. Nam
sollicitudin elit id risus tempor congue. Cras fringilla, libero a faucibus feugiat, odio lacus
vulputate ligula,
vel euismod tellus orci eget lorem \cite{citation11, citation12,citation13, citation14,citation15,
citation16, citation17, citation18, citation19, citation20}.
\section{References}
\begin{enumerate}
\item
T. He, H. Khamfroush, S. Wang, T. La Porta, and S. Stein, “It’s hard to
share: Joint service placement and request scheduling in edge clouds with
sharable and non-sharable resources,” in IEEE 38th International Conference on Distributed
Computing Systems (ICDCS), 2018, pp. 365–375.
\item
K. Poularakis, J. Llorca, A. M. Tulino, I. Taylor, and L. Tassiulas, “Joint
service placement and request routing in multi-cell mobile edge computing networks,” in
IEEE Conference on Computer Communications (INFOCOM), 2019, pp. 10–18.
\item
K. Poularakis, J. Llorca, A. M. Tulino, and I. Taylor, “Service placement
and request routing in mec networks with storage, computation, and communication
constraints,” IEEE/ACM Transactions on Networking, vol. 28,
no. 3, pp. 1047–1060, 2020.
\item
V. Farhadi, F. Mehmeti, T. He, T. L. Porta, H. Khamfroush, S. Wang, and
K. S. Chan, “Service placement and request scheduling for data-intensive 1
applications in edge clouds,” in IEEE Conference on Computer
Communications(INFOCOM), 2019, pp. 1279–1287.
\item
Y. Qu, H. Dai, H. Wang, C. Dong, F. Wu, S. Guo, and Q. Wu, “Service
provisioning for uav-enabled mobile edge computing,” IEEE Journal on
Selected Areas in Communications, vol. 39, no. 11, pp. 3287–3305, 2021.
\item
X. Xu, H. Cao, Q. Geng, X. Liu, F. Dai, and C. Wang, “Dynamic resource
provisioning for workflow scheduling under uncertainty in edge computing
environment,” Concurrency and Computation: Practice and Experience,
vol. n/a, no. n/a, p. e5674.
\item
B. Li, Q. He, G. Cui, X. Xia, F. Chen, H. Jin, and Y. Yang, “Read:
Robustness-oriented edge application deployment in edge computing environment,” IEEE
Transactions on Services Computing, pp. 1–1, 2020.
\item
NSU, “Rsome,” https://xiongpengnus.github.io/rsome/, 2021, [Online; accessed 3-Feb-2022].
\item
N. Ngoenriang, S. J. Turner, D. Niyato, and S. Supittayapornpong, “Joint
uav-placement and data delivery in aerial inspection under uncertainties,”
IEEE Internet of Things Journal, pp. 1–1, 2021.
\item
IBM, “Ibm cplex optimizer,” https://www.ibm.com/in-en/analytics/cplexoptimizer, 2021,
[Online; accessed 3-Feb-2022].
\end{enumerate}
\end{document}
